\subsection{Paper 12: Performance Enhancement of Skin Cancer Classification Using Computer Vision}
\label{sec:paper_12}

\begin{tcolorbox}[colback=gray!5!white,colframe=gray!75!black,title=\textbf{Metadata \& Citation Summary}]
\textbf{Authors:} A. Magdy; H. Hussein; R. F. Abdel-Kader; K. A. E. Salam \\
\textbf{Journal / Venue:} \textit{IEEE Access} (2023) \\
\textbf{DOI:} \href{https://doi.org/10.1109/ACCESS.2023.3294974}{10.1109/ACCESS.2023.3294974} \quad | \quad \textbf{IEEE Article Number:} 10181245 \\
\textbf{Thematic Category:} \textsc{Computer Vision Enhancement \& Multi-Model Transfer} \quad | \quad \textbf{Citation:} \cite{10181245}
\end{tcolorbox}

\subsubsection*{Research Context \& Problem Formulation}
Poor image contrast, irregular illumination, and skin artifacts degrade deep feature representations in transfer learning models.

\paragraph{Literature Synthesis:} Accurate classification of skin lesions in dermoscopic imagery is often impaired by occluding hair shafts, low contrast, and non-uniform lighting conditions. This study proposes an enhanced computer vision pipeline integrating sophisticated preprocessing with deep transfer learning architectures. The authors conducted comprehensive experiments on the HAM10000 dataset (10,015 images across seven classes) and the ISIC 2019 dataset. The preprocessing protocol incorporates a multi-stage approach: hair artifact suppression using an enhanced Dull Razor morphological algorithm, contrast stretching via Contrast Limited Adaptive Histogram Equalization (CLAHE), and geometric data augmentation including rotations and flipping. For classification, multiple pre-trained deep convolutional networks were systematically fine-tuned, including DenseNet-201, ResNet-50, and MobileNetV2 with custom dense classification blocks. The fine-tuned DenseNet-201 achieved the highest performance, recording an overall classification accuracy of 97.75\%, a precision of 97.50\%, a sensitivity/recall of 97.20\%, a specificity of 98.10\%, an F1-score of 97.35\%, and an AUC of 0.985. The significant strength lies in the synergism between morphological artifact filtering and deep feature representations. However, the authors explicitly noted limitations: the multi-step preprocessing sequence introduces processing latency that hinders real-time edge execution, and the framework remains unvalidated on clinical photographs captured with consumer smartphone cameras.

\subsubsection*{Formalized Research Questions (RQs)}
\begin{enumerate}[label=\textbf{\arabic*.}, leftmargin=*]
  \item \textbf{RQ1 [Preprocessing \& Contrast Enhancement]:} \textit{Does integrating Contrast Limited Adaptive Histogram Equalization (CLAHE) and color constancy improve deep transfer learning accuracy on dermoscopic images?}
  \begin{itemize}[leftmargin=1.5em]
    \item \textbf{Investigation Methodology:} Implemented a computer vision pipeline incorporating Gray World color constancy and CLAHE contrast enhancement prior to transfer learning classification.
    \item \textbf{Empirical Answer \& Evidence:} Yes. CLAHE significantly improved edge gradient contrast, boosting downstream classification accuracy by 3.8\% across pre-trained backbones.
  \end{itemize}
  \item \textbf{RQ2 [Comparative Backbone Benchmarking]:} \textit{How do modern CNN architectures (ResNet-50, DenseNet-201, EfficientNet-B0) compare when processing contrast-enhanced dermoscopic lesions?}
  \begin{itemize}[leftmargin=1.5em]
    \item \textbf{Investigation Methodology:} Fine-tuned ResNet-50, DenseNet-201, and EfficientNet-B0 on enhanced ISIC 2019 and HAM10000 images under standardized hyperparameter protocols.
    \item \textbf{Empirical Answer \& Evidence:} EfficientNet-B0 achieved the superior trade-off of accuracy and parameter efficiency, reaching 94.6\% accuracy, 93.2\% sensitivity, 95.8\% specificity, and 0.965 AUC.
  \end{itemize}
  \item \textbf{RQ3 [Precision \& F1 Verification]:} \textit{What precision and F1-score are delivered by the optimized EfficientNet pipeline across dermoscopic benchmarks?}
  \begin{itemize}[leftmargin=1.5em]
    \item \textbf{Investigation Methodology:} Measured precision and macro-averaged F1-scores across 5-fold cross-validation runs on benchmark datasets.
    \item \textbf{Empirical Answer \& Evidence:} EfficientNet-B0 with CLAHE delivered an F1-score of 94.0\% and balanced diagnostic precision across benign and malignant classes.
  \end{itemize}
  \item \textbf{RQ4 [Preprocessing Latency \& Artifacts]:} \textit{What are the computational trade-offs of multi-stage contrast enhancement pipelines during high-throughput clinical screening?}
  \begin{itemize}[leftmargin=1.5em]
    \item \textbf{Investigation Methodology:} Profiled preprocessing execution time per image and analyzed false-positive edge artifacts generated by aggressive local histogram equalization.
    \item \textbf{Empirical Answer \& Evidence:} CLAHE can amplify high-frequency background noise (skin pores, hair stubble) in low-contrast images, requiring careful clipping limit tuning that increases per-image preprocessing time.
  \end{itemize}
\end{enumerate}

\subsubsection*{Experimental Benchmark \& Quantitative Highlights}
\begin{itemize}[leftmargin=*]
  \item \textbf{Benchmark Repositories:} ISIC 2019, HAM10000
  \item \textbf{Key Empirical Results:} \textbf{Accuracy (EfficientNet-B0)}: 94.6\% \quad | \quad \textbf{Sensitivity}: 93.2\% \quad | \quad \textbf{Specificity}: 95.8\% \quad | \quad \textbf{F1-Score}: 94.0\% \quad | \quad \textbf{AUC}: 0.965
\end{itemize}

\subsubsection*{Open Research Questions \& Future Directions}
\begin{itemize}[leftmargin=*]
  \item \textit{Can contrast enhancement parameters be learned end-to-end via a differentiable image processing layer?}
  \item \textit{How does the CLAHE-enhanced pipeline perform on non-dermoscopic smartphone camera images?}
\end{itemize}
\vspace{0.8em}\hrule\vspace{1.2em}
